\documentclass[10pt,a4paper]{amsart}
\usepackage[margin=19mm]{geometry}
\usepackage{amsmath,amssymb,mathtools}
\usepackage[scr=rsfs,cal=euler]{mathalfa}
\usepackage{xfrac}
\usepackage[unicode,hidelinks,bookmarks=false]{hyperref}
\newcommand{\ie}{i.e.\ }
\newcommand{\cf}{cf.\ }
\newcommand{\resp}{respectively\ }
\newcommand{\Subsection}{\subsection}
\providecommand{\nobreakdash}{\nobreak\hbox{-}}
\setlength{\emergencystretch}{2em}
\multlinegap0pt
\usepackage{mathtools}
\usepackage{bm}
\usepackage{dsfont}
\usepackage{xcolor}
\usepackage{cancel}
\usepackage{amssymb}
\usepackage{tikz}
\newtheorem{definition}{Definition}
\newtheorem{theorem}{Theorem}
\newtheorem{corollary}{Corollary}
\newtheorem{lemma}{Lemma}
\newtheorem{remark}{Remark}
\newcommand{\RR}{{\mathbb R}}
\newcommand{\ZZ}{{\mathbb Z}}
\newcommand{\id}{{\operatorname{id}}}
\newcommand{\tZ}{{\widetilde{\mathbb Z}}}
\title{Statistics of Abelian topological excitations}
\author{Hanyu Xue}
\date{Source: arXiv:2412.07653v5 (CC BY 4.0)}
\begin{document}
\maketitle

\noindent\emph{Source and licence.} Statistics of Abelian topological excitations. Version arXiv:2412.07653v5 (CC BY 4.0), distributed by arXiv under the Creative Commons Attribution 4.0 International licence, http://creativecommons.org/licenses/by/4.0/, as recorded in the arXiv licence metadata for this exact version and shown on the abstract page https://arxiv.org/abs/2412.07653v5. Full licence notice: \texttt{licenses/arxiv-2412.07653-CC-BY-4.0.txt}.\par\smallskip

\noindent\textit{Editorial scope.} Counted unit: the article's lemma lemma (source0.tex lines 2760--2762). The article's complete original proof of it is delivered with it (source0.tex lines 2763--2822, one \texttt{proof} environment of 60 lines). No mathematics is written, repaired or re-derived: the statement and the proof are the article's own lines, in the article's own notation, and no retained line is substituted or re-flowed.  Retained uncounted as the source-local context the proof needs: lines 339--347; lines 556--560; lines 629--634; lines 2478--2480; lines 2604--2626; lines 2730--2732; lines 2743--2751.  Nothing was omitted from any retained span, so the package carries no omission record.  No retained line contains a citation, so the wrapper carries no bibliography and no bibliography entry.  No retained line contains a figure environment, an image-inclusion command or drawing code, so no image is delivered and the package declares no asset dependency.  Editorial formatting only, in addition: an AMS \texttt{amsart} preamble that restates the article's own declarations and macros used by the retained lines, namely the environments \texttt{definition}, \texttt{theorem}, \texttt{corollary}, \texttt{lemma}, \texttt{remark} on separate counters, together with the standard packages those lines need.  The retained environments print in this document as Definition 1, Theorem 1, Corollary 1, Lemma 1, Remark 1 in order of appearance, whereas the article prints the same items with the numbering of its own full document.  The complete native source of the article as distributed in the arXiv e-print for this exact version is bundled unchanged as \texttt{sources/169903/source0.tex}.
\begin{definition}\label{expSimplicialComplex}
	The excitation pattern $m_p(X,G)$ is constructed using the following data:
	\begin{enumerate}
		\item \(A\) is the group of \(p\)-dimensional simplicial boundaries, \(A = B_p(X, G)\).
		\item \(S = G_0\times X_{p+1}\), where \(X_{p+1}\) is the set of all \((p+1)\)-simplexes of $X$, and \(G_0\) is a generating subset of \(G\). 
		\item Viewing $X$ as a topological space, \(\operatorname{supp}(s)\) gives the geometric image of the simplex itself, while \(\partial\) is the usual homological boundary map.
	\end{enumerate}

\end{definition}

\begin{theorem}\label{conjectureCohomology}
	\begin{equation}
		T^*(m_p(\partial\Delta^{d+1},G))\simeq H^{d+2}(K(G,d-p),\RR/\ZZ).
	\end{equation}
\end{theorem}

\begin{equation}\label{eqTheta(g,a)}
	\begin{aligned}
		(g_2g_1,a)&=(g_1,a)+(g_2,a+\partial g_1),\\
		(g^{-1},a)&=-(g,a-\partial g).
	\end{aligned}
\end{equation}

\begin{corollary}\label{corollaryZeroHomotopic}
	Let $ L\xrightarrow{p} K$ makes $L$ a path space of $K$, then a spherical $n$-simplex $\sigma \in K_n$ is zero-homotopic if and only if there is a spherical $n$-simplex $\tau\in L_n$ satisfying $p_n(\tau)=\sigma$.
\end{corollary}

\begin{lemma}\label{lemmaCochainTool}
	For any $0\le k\le n$, we define a linear map $\iota_k:C^{q-1}(\Delta^{n-1},G)\rightarrow C^{q-1}(\Delta^{n},G)$ by identifying $\Delta^{n-1}$ with $\partial_k\Delta^n$ and extending to $\Delta^n$ by zero. Then, we have
	\begin{enumerate}
		\item The composition of 
		\begin{equation}
			C^{q-1}(\Delta^{n-1},G)\xrightarrow{\iota_k}C^{q-1}(\Delta^{n},G)\xrightarrow{d}Z^{q}(\Delta^{n},G)
		\end{equation}
		is an isomorphism $C^{q-1}(\Delta^{n-1},G)\simeq Z^{q}(\Delta^{n},G)$.
		\item 
		\begin{equation}
			\partial_i\iota_k=\left\{
			\begin{aligned}
				\id\quad \quad & \text{if } i=k;\\
				\iota_{k-1}\partial_i\quad & \text{   if } i<k;\\
				\iota_{k}\partial_{i-1}\quad  & \text{   if } i>k.
			\end{aligned}
			\right.
		\end{equation}
		Here, $\partial_i$ is restricting a cochain on the $i$-th face.
		
		
	\end{enumerate}
\end{lemma}

 \begin{equation}\label{eqB37}
 	\mathcal{D}([\sum_{i=1}^N \alpha_ix_i,\beta])=[0,\beta]+\sum_{i=1}^N\bigg([\alpha_ix_i,d(\sum_{j=1}^{i-1}\alpha_jx_j)+\beta]-[0,d(\sum_{j=1}^{i-1}\alpha_jx_j)+\beta]\bigg)
 \end{equation}

 \begin{remark}\label{remarkDecomposition}
 	The intuition of $\mathcal{D}$ is similar to Eq.~\eqref{eqTheta(g,a)}, where we decompose the action of a process $(g,a)$ to elementary steps $(s_i,a_i)$. In this section, we view $[\sum_{i=1}^N\alpha_ix_i,\beta]$ as a successive action of $\alpha_1x_1,\alpha_2x_2,\cdots$; in the first step, the action of $\alpha_1x_1$ produce a term $[\alpha_1x_1,\beta]$, and the remain process corresponds to $[\sum_{i=2}^N\alpha_ix_i,d(\alpha_1x_1)+\beta]$. Thus, we get the decomposition
 	\begin{equation}
 		[\sum_{i=1}^N \alpha_ix_i,\beta]\mapsto [\alpha_1x_1,\beta]+[\sum_{i=2}^N\alpha_ix_i,d(\alpha_1x_1)+\beta]-[0,d(\alpha_1x_1)+\beta],
 	\end{equation}
 	where the last term is used to compensate $d_0$ and $d_1$. Keep doing similar decomposition, one will get Eq.~\eqref{eqB37} in the end.
 	
 	One may raise the question that why in Definition~\ref{expSimplicialComplex}, we use $S=G_0\times X_{p+1}$ but not $S=C_{p+1}(X,G)$ to define $m_p(X,G)$. The reason is from the locality axiom: a simplex naturally has a support, while the support of a chain is a bad notion. Actually, we have tried to use $S=C_{p+1}(X,G)$ as the set of excitation operators and to replace the locality axiom by some other conditions, but none of them is physically basic and produce Theorem~\ref{conjectureCohomology} at the same time.
 \end{remark}

\begin{lemma}
	If $\lambda\in \mathcal{Z}$, then $\lambda-\mathcal{D}(\lambda)\in\mathcal{B}$.
\end{lemma}

\begin{proof}
	Following Remark~\ref{remarkDecomposition}, without loss of generality, we prove that for any
	\begin{equation}
		\lambda=\sum_{\alpha\in C^{q-1}(\Delta^{d+1},G),\beta\in Z^q(\Delta^{d+1},G)} c(\alpha,\beta)[\sum_{i=1}^N \alpha_ix_i,\beta]\in \mathcal{Z},
	\end{equation}
	we have
	\begin{equation}
		\mu=\sum c(\alpha,\beta)\bigg([\sum_{i=1}^N \alpha_ix_i,\beta]-[\alpha_1x_1,\beta]-[\sum_{i=2}^N \alpha_ix_i,\alpha_1dx_1+\beta]+[0,\alpha_1dx_1+\beta]\bigg)\in \mathcal{B}.
	\end{equation}
	The idea is to use Corollary~\ref{corollaryZeroHomotopic}. We have already constructed a path fibration sequence $\Omega\tZ[K(G,q)]\rightarrow L\rightarrow \tZ[K(G,q)]$; in parallel, we construct another path fibration sequence $\Omega\tZ[L(G,q)]\rightarrow L'\rightarrow \tZ[L(G,q)]$. More explicitly, the $\Omega\tZ[L(G,q)]$ is the subgroup $\ker d_0\cap \ker d_1$ of $\ZZ[L(G,q)\times L(G,q)]$, where 
	$d_0,d_1:\ZZ[L(G,q)\times K(G,q)]\rightarrow \ZZ[K(G,q)]$ are defined by
	\begin{equation}
		\left\{
		\begin{aligned}
			d_0([\alpha,\gamma]) &= [\alpha],\\
			d_1([\alpha,\gamma]) &= [\alpha+\gamma],
		\end{aligned}
		\right.
	\end{equation}
	where $\alpha,\gamma\in C^{q-1}(\Delta^n,G)$. There is a canonical simplicial map
	\begin{equation}
		\begin{aligned}
			\delta : \Omega\tZ[L(G,q)]&\rightarrow \Omega\tZ[K(G,q)]\\
			[\alpha,\gamma]&\mapsto[\alpha,d\gamma].
		\end{aligned}		
	\end{equation}
	
	We fix an integer $k_0$ such that $0\le k_0\le d+1$ and $x_1\nsubseteq \partial_{k_0}\Delta^{d+1}$. Using Lemma~\ref{lemmaCochainTool}, we have an isomorphism $d\circ \iota_{k_0}:C^{q-1}(\Delta^{d},G)\rightarrow Z^q(\Delta^{d+1},G)$. Replacing $\beta$ by $d \iota_{k_0}\gamma$ where $\gamma\in C^{q-1}(\Delta^{d},G)$, we have
	\begin{equation}
		\lambda=\sum_{\alpha,\gamma} c(\alpha,d \iota_{k_0}\gamma)[\sum_{i=1}^N \alpha_ix_i,d \iota_{k_0}\gamma]
	\end{equation}
	and
	\begin{equation}
		\mu=\sum_{\alpha,\gamma} c(\alpha,d \iota_{k_0}\gamma)\bigg([\sum_{i=1}^N \alpha_ix_i,d \iota_{k_0}\gamma]-[\alpha_1x_1,d \iota_{k_0}\gamma]-[\sum_{i=2}^N \alpha_ix_i,\alpha_1dx_1+d \iota_{k_0}\gamma]+[0,\alpha_1dx_1+d \iota_{k_0}\gamma]\bigg)
	\end{equation}
	Next, we define
	\begin{equation}\label{eqB54}
		\rho=\sum_{\alpha,\gamma} c(\alpha, d\iota_{k_0}\gamma)\bigg([\sum_{i=1}^N \alpha_ix_i, \iota_{k_0}\gamma]-[\alpha_1x_1, \iota_{k_0}\gamma]-[\sum_{i=2}^N \alpha_ix_i,\alpha_1x_1+ \iota_{k_0}\gamma]+[0,\alpha_1x_1+ \iota_{k_0}\gamma]\bigg)
	\end{equation}
	
	It is obvious that $d\rho=\mu$ and $d_0\rho=d_1\rho=0$. We now show that $\partial_k\rho=0$ for all $0\le k\le d+1$.
	
	
	If $x_1\nsubseteq \partial_{k}\Delta^{d+1}$, then in $\partial_k\mu$, all $\alpha_1x_1$ is replaced by zero, and the result is automatically zero.
	
	If $x_1\subseteq \partial_{k}\Delta^{d+1}$, we have $k\ne k_0$ by the selection of $k_0$. Because $\lambda\in\mathcal{Z}$, we have
	
	\begin{equation}
		\partial_k\lambda=\sum_{\alpha,\gamma} c(\alpha,d \iota_{k_0}\gamma)\bigg[\sum_{i=1}^N\partial_k \alpha_ix_i,d\partial_k\iota_{k_0}\gamma\bigg]=0.
	\end{equation}
	Using Lemma~\ref{lemmaCochainTool} (the case for $k<k_0$ and $k>k_0$ is similar, we assume $k>k_0$), we get
	\begin{equation}
		\sum_{\alpha,\gamma} c(\alpha,d \iota_{k_0}\gamma)\bigg[\sum_{i=1}^N\partial_k \alpha_ix_i,\partial_{k-1}\gamma\bigg]=0.
	\end{equation}
	
	This equation is strong enough to imply that all of the four terms in Eq.~\eqref{eqB54} is zero. In all, we have $\rho\in \cap_{i=0}^{d+1}\ker \partial_i$.
	
	Finally, we use $d\circ \iota_{d+2}$ to map $\rho$ into $\Omega\tZ[K(G,q)]_{d+2}$. Using Lemma~\ref{lemmaCochainTool}, we have $\partial_{d+2}\circ d\circ \iota_{d+2}(\rho)=\mu$ and $\partial_{i}\circ d\circ \iota_{d+2}(\rho)=0$ for $i<d+2$.
	This finishes the proof of this lemma and also Theorem~\ref{conjectureCohomology}.
\end{proof}

\end{document}
